\section{Principal Symbol and Hyperbolicity}

To analyze well-posedness we construct the principal symbol of the coupled Einstein--Proca system.

The state vector is:

\begin{equation}
\Psi = (g_{ij},K_{ij},A_\mu,\phi)
\end{equation}

The principal symbol is defined through the highest derivative terms:

\begin{equation}
\mathbb P(\xi) \Psi = 0
\end{equation}

where \(\xi_\mu\) is the characteristic covector.

Strong hyperbolicity requires:

\begin{itemize}
\item real characteristic speeds
\item complete eigenvector basis
\item bounded condition number
\end{itemize}

The longitudinal Proca mode yields the dominant stability constraint.

In the WKB limit:

\begin{equation}
v_L^2
\approx
1
-
\frac{\alpha^2 \dot\phi^2}{M_*^2 M_A^2}
\end{equation}

Therefore:

\begin{equation}
\alpha < \alpha_{max}
\end{equation}

is required to avoid Laplacian instability.

The perturbative stability island is:

\begin{equation}
\alpha_c < \alpha < \alpha_{max}
\end{equation}

Outside this region the PDE system becomes ill-posed.
